\documentclass[a4paper,12pt]{article}

\usepackage[italian]{babel}
\usepackage[T1]{fontenc}
\usepackage[utf8]{inputenc}
\usepackage{geometry}
\usepackage{hyperref}
\usepackage{graphicx}
\usepackage{longtable}
\usepackage{setspace}

\geometry{margin=2.5cm}
\onehalfspacing

\title{\textbf{DOCUMENTAZIONE TECNICA DI NUTRIAPP}}
\author{Team NutriApp}
\date{\today}

\begin{document}

    \maketitle
    \tableofcontents
    \newpage

    \section{Introduzione}

    NutriApp è un'applicazione mobile progettata per il monitoraggio intelligente di calorie, macronutrienti e micronutrienti.
    L'obiettivo principale del progetto è fornire agli utenti uno strumento digitale affidabile e intuitivo per sviluppare maggiore consapevolezza alimentare e supportare il raggiungimento di obiettivi nutrizionali personalizzati.

    L'applicazione consente la registrazione dei pasti tramite tre modalità principali: inserimento manuale degli ingredienti, utilizzo di ricette salvate e scansione del codice a barre dei prodotti confezionati. I dati raccolti vengono elaborati per fornire un'analisi quantitativa dell'apporto nutrizionale e una visualizzazione storica dell'andamento nel tempo.

    \section{Che problemi risolve}

    Uno dei principali ostacoli al mantenimento di uno stile di vita sano è la mancanza di feedback strutturato e continuo. Molti utenti non dispongono di strumenti adeguati per monitorare con precisione l'assunzione di nutrienti, rendendo difficile valutare l'efficacia delle proprie scelte alimentari.

    NutriApp affronta questa criticità offrendo un sistema di monitoraggio costante, basato su dati aggiornati e rappresentazioni grafiche intuitive. L'utente può così modificare le proprie abitudini alimentari in funzione dei risultati osservati, riducendo errori di valutazione e aumentando la motivazione personale.

    \subsection{La richiesta del committente}

    La committente ha espresso la necessità di sviluppare un'applicazione altamente personalizzabile, in grado di consentire agli utenti di selezionare i parametri nutrizionali da monitorare con priorità.

    In risposta a tale richiesta, l'architettura dell'app è stata progettata per permettere la configurazione dinamica degli obiettivi nutrizionali e dei nutrienti monitorati, garantendo flessibilità e adattabilità alle diverse esigenze individuali.

    \subsection{L'interesse delle masse}

    L'interesse generale verso uno stile di vita sano e consapevole è in costante crescita. Parallelamente, aumenta la necessità di strumenti digitali in grado di supportare persone con condizioni specifiche quali diabete, ipertensione o intolleranze alimentari.

    NutriApp è stata progettata con un approccio inclusivo, mirando a supportare sia utenti generici interessati al miglioramento dello stile alimentare sia soggetti con esigenze nutrizionali particolari.

    \section{Lo sviluppo dell'applicazione}

    Lo sviluppo del progetto è stato condotto adottando la metodologia Agile. Il lavoro è stato suddiviso in sprint incrementali, ciascuno focalizzato sull'implementazione di funzionalità specifiche.

    Questa scelta metodologica ha permesso un monitoraggio continuo dell'avanzamento, una rapida identificazione di criticità e una progressiva validazione delle componenti sviluppate.

    \subsection{Analisi dei requisiti e degli stakeholders}

    \subsubsection{Requisiti funzionali e tecnici}

    I requisiti principali individuati in fase di analisi includono:

    \begin{itemize}
        \item Inserimento manuale dei pasti tramite elenco ingredienti e quantità.
        \item Registrazione di ricette personalizzate riutilizzabili.
        \item Inserimento tramite scansione di codice a barre.
        \item Calcolo automatico di calorie, macronutrienti e micronutrienti.
        \item Visualizzazione grafica dei dati storici su base giornaliera, settimanale e mensile.
        \item Sistema di notifiche per il promemoria della registrazione pasti.
    \end{itemize}

    L'applicazione integra l'API OpenFoodFacts per il recupero dei dati nutrizionali. I valori vengono adattati alle quantità inserite dall'utente e aggregati nel diario personale.

    Il sistema è progettato per garantire accuratezza nei calcoli e coerenza nei dati visualizzati.

    il documento per l'analisi dei requisiti dell'applicazione è disponibile a questo collegamento:

    \href{https://gitlab.com/ismailbarakat68/nutriapp/-/raw/documentazioneSprint3/Documentazione/documentazioneLatex/analisi_requisiti.pdf?ref_type=heads&inline=true}{LINK}


    \subsubsection{Stakeholders}

    Gli stakeholder identificati comprendono:

    \begin{itemize}
        \item La committente, con ruolo decisionale e strategico.
        \item Gli utenti finali con esigenze nutrizionali specifiche.
        \item Gli utenti orientati al miglioramento dello stile di vita.
        \item Il fornitore dell'API nutrizionale, da cui dipende la qualità dei dati.
        \item Professionisti della nutrizione potenzialmente interessati all'utilizzo dello strumento.
        \item Strutture ospedaliere e sanitarie.
    \end{itemize}

    L'analisi degli stakeholder ha permesso di identificare rischi, dipendenze e priorità progettuali.

    Il documento di analisi degli stakeholder è visibile a questo collegamento:

    \href{https://gitlab.com/ismailbarakat68/nutriapp/-/raw/documentazioneSprint3/Documentazione/documentazioneLatex/analisi_stakeholders.pdf?ref_type=heads&inline=true}{LINK}


    \subsection{Creazione del Project Charter}

    il project charter è visibile a questo link:

    \href{https://gitlab.com/ismailbarakat68/nutriapp/-/raw/documentazioneSprint3/Documentazione/documentazioneLatex/project_charter.pdf?ref_type=heads&inline=true}{LINK}


    Il Project Charter rappresenta il documento formale di avvio del progetto NutriApp.
    Esso definisce obiettivi, ambito, stakeholder principali, milestone temporali, budget e vincoli operativi.
    Il Project Charter individua anche i principali rischi:
    dipendenza da API esterne (OpenFoodFacts), possibili modifiche delle politiche degli store digitali e limitata disponibilità del team di sviluppo.

    \subsection{Divisione del lavoro}

    La suddivisione delle attività è stata organizzata secondo la metodologia Agile, attraverso sprint incrementali basati su un Product Backlog strutturato in User Stories.

    Le User Stories sono state formulate secondo il modello:
    `Come [ruolo] vorrei [azione] per [scopo]''

    Le funzionalità sono state classificate per priorità:

    Priorità 1 (MUST):
    interfaccia intuitiva, inserimento manuale, scansione barcode, calcolo automatico dei totali giornalieri e sistema di notifiche push.

    Priorità 2 (SHOULD):
    registrazione account, salvataggio ricette personalizzate, impostazione soglie nutrizionali personalizzate e generazione grafici giornalieri.

    Priorità 3 (MAY):
    report periodici settimanali/mensili e tracciamento di parametri fisici.

    Le tempistiche stimate per ciascuna funzionalità variano tra 6 e 18 giorni lavorativi, consentendo una pianificazione progressiva e modulare dello sviluppo.

    Questa organizzazione ha permesso di garantire iterazioni rapide, validazione continua e adattamento dinamico alle esigenze emergenti.

    \subsection{Divisione dei ruoli e avvio dello sviluppo}

    L'organizzazione del team è stata definita in base alle competenze specifiche dei membri:

    La Prof.ssa Manuela Mileo ricopre il ruolo di committente, con alto livello di interesse e potere decisionale, fornendo linee guida strategiche e supervisione.

    Barakat Ismail svolge il ruolo di Project Manager e responsabile Backend e Data Engineering, coordinando le attività tecniche e l'integrazione con i servizi esterni.

    Maltese Emanuel è responsabile dello sviluppo Frontend.

    Yu Serena cura la progettazione UX/UI, assicurando coerenza grafica e usabilità.

    Donzelli Christian opera come esperto nutrizionale,ricercatore delle informazioni scientifiche e scrittore della documentazione inerente al progetto.

    Gli stakeholder esterni includono utenti con necessità nutrizionali specifiche, utenti generici interessati al miglioramento dello stile alimentare, professionisti della nutrizione, strutture ospedaliere e il fornitore del database nutrizionale.

    Lo sviluppo è stato avviato dopo la definizione dell’architettura tecnica e l’approvazione del design UI/UX, con una pianificazione temporale articolata in milestone che culminano nel rilascio previsto entro Giugno 2026.

    \section{Guida all'utilizzo dell'applicazione}

    link del documento:
    \href{https://gitlab.com/ismailbarakat68/nutriapp/-/raw/2db73b6a2bec3960f3ce16a5863640902789d547/Documentazione/documentazioneLatex/guida_di_utilizzo.pdf?inline=true}{Apri il documento di guida all'utilizzo}


    \section{Fonti dei dati e ricerca scientifica}

    La presente sezione documenta le analisi nutrizionali e scientifiche svolte in fase preliminare allo sviluppo dell’applicazione NutriApp.
    Le valutazioni effettuate hanno riguardato la selezione delle fonti dati, l’identificazione dei nutrienti da monitorare, l’analisi delle principali condizioni cliniche correlate all’alimentazione e lo studio delle formule per la determinazione del fabbisogno energetico individuale.

    \subsection{Integrazione dell'API OpenFoodFacts}

    Per l’acquisizione dei dati nutrizionali relativi ai prodotti alimentari con codice a barre è stata selezionata l’API OpenFoodFacts, un database open source contenente informazioni su alimenti confezionati e non confezionati distribuiti a livello internazionale.

    L’API consente il recupero dei valori nutrizionali standardizzati per 100 grammi di prodotto, includendo energia, macronutrienti, micronutrienti, ingredienti e allergeni. L’integrazione avviene mediante interrogazione basata su codice a barre (EAN), con successiva normalizzazione dei dati ricevuti in formato JSON.

    Durante la fase di analisi tecnica sono state riscontrate possibili criticità relative alla presenza di campi mancanti, formati non uniformi e incompletezza dei micronutrienti. È stato pertanto progettato un sistema di validazione, gestione delle eccezioni e fallback dei valori non disponibili, al fine di garantire coerenza nei calcoli nutrizionali.

    \subsection{Macronutrienti e parametri energetici}

    I macronutrienti costituiscono il nucleo centrale del sistema di monitoraggio dell’applicazione.

    I carboidrati rappresentano la principale fonte energetica dell’organismo e sono fondamentali per il funzionamento del sistema nervoso centrale. Essi vengono convertiti in glucosio e contribuiscono al mantenimento della glicemia. Il monitoraggio include carboidrati totali, zuccheri semplici e fibre. Particolare attenzione è stata dedicata alla distribuzione giornaliera, all’indice glicemico e al carico glicemico, in quanto un’assunzione eccessiva o mal distribuita è associata a iperglicemia, insulino-resistenza e sindrome metabolica, mentre una carenza può determinare affaticamento e chetosi.

    Le proteine, costituite da aminoacidi essenziali e non essenziali, svolgono funzioni strutturali, enzimatiche, immunitarie e ormonali. Il monitoraggio avviene in grammi per chilogrammo di peso corporeo, considerando la distribuzione nell’arco della giornata e l’origine animale o vegetale. Carenze proteiche possono determinare sarcopenia e immunodepressione, mentre eccessi risultano potenzialmente critici nei soggetti con insufficienza renale.

    I grassi rappresentano una fonte energetica concentrata e svolgono un ruolo essenziale nella struttura delle membrane cellulari e nell’assorbimento delle vitamine liposolubili. L’applicazione distingue tra grassi saturi, monoinsaturi, polinsaturi (omega-3 e omega-6) e trans. È stata prevista un’analisi qualitativa del rapporto tra grassi saturi e insaturi, data la correlazione con dislipidemia e rischio cardiovascolare.

    L’energia totale viene calcolata in chilocalorie (kcal) sulla base dei valori nutrizionali forniti dall’API e delle quantità inserite dall’utente.

    \subsection{Calcolo del fabbisogno energetico}

    Per la determinazione del metabolismo basale (BMR) è stata adottata la formula di Mifflin-St Jeor, considerata tra le più accurate in ambito clinico.

    Per gli uomini:
    \[
        BMR = 10 \times peso (kg) + 6.25 \times altezza (cm) - 5 \times età (anni) + 5
    \]

    Per le donne:
    \[
        BMR = 10 \times peso (kg) + 6.25 \times altezza (cm) - 5 \times età (anni) - 161
    \]

    Il BMR viene successivamente moltiplicato per un fattore di attività fisica (FA) al fine di ottenere il TDEE (Total Daily Energy Expenditure). I fattori utilizzati sono:

    1.2 per soggetti sedentari
    1.375 per soggetti leggermente attivi
    1.55 per soggetti moderatamente attivi
    1.725 per soggetti molto attivi
    1.9 per soggetti con attività intensa o lavoro fisico pesante

    Questo calcolo consente all’applicazione di personalizzare gli obiettivi calorici giornalieri in funzione delle caratteristiche individuali.

    \subsection{Micronutrienti: vitamine e minerali}

    Le vitamine sono state suddivise in liposolubili (A, D, E, K) e idrosolubili (complesso B e vitamina C).
    Sono stati analizzati rischi di carenza ed eccesso, nonché popolazioni a rischio specifiche. Ad esempio, la vitamina D è frequentemente carente nella popolazione generale ed è associata alla salute ossea; la vitamina B12 è critica nei soggetti vegani e negli anziani; i folati risultano fondamentali in gravidanza.

    Per quanto riguarda i minerali, sono stati distinti macro-minerali (calcio, fosforo, magnesio, sodio, potassio) e oligoelementi (ferro, zinco, iodio, selenio).
    Particolare attenzione è stata posta al sodio per la correlazione con ipertensione, al ferro per il rischio di anemia, al calcio e vitamina D per osteoporosi, e al potassio per alterazioni elettrolitiche.

    \subsection{Fibre, acqua ed elettroliti}

    Le fibre alimentari, suddivise in solubili e insolubili, svolgono un ruolo determinante nella regolazione della funzione intestinale, nel controllo della glicemia postprandiale e nella modulazione del colesterolo. Un apporto insufficiente è associato a stipsi, disbiosi e aumentato rischio metabolico.

    L’acqua è essenziale per il mantenimento dell’omeostasi, la termoregolazione e il trasporto dei nutrienti. Gli elettroliti contribuiscono all’equilibrio dei fluidi e alla funzione neuromuscolare. Popolazioni a rischio di disidratazione includono anziani, sportivi e pazienti con patologie renali.

    \subsection{Supporto per condizioni cliniche specifiche}

    L’analisi preliminare ha individuato diverse condizioni cliniche in cui il monitoraggio nutrizionale risulta particolarmente rilevante.

    Nel diabete è prioritario il controllo dei carboidrati e delle fibre.
    Nell’ipertensione è centrale il monitoraggio del sodio e del potassio.
    Nell’osteoporosi assumono rilevanza calcio e vitamina D.
    Nell’anemia risultano critici ferro, vitamina B12 e folati.
    Nella celiachia è necessario monitorare ferro, B12 e fibre.
    In gravidanza sono fondamentali folati, iodio e ferro.
    Nel veganismo occorre attenzione a vitamina B12, ferro e omega-3.
    Nell’insufficienza renale è necessario controllare proteine, potassio e fosforo.

    L’applicazione consente l’impostazione di soglie personalizzate per nutrienti specifici, con sistema di monitoraggio continuo e segnalazione del superamento dei limiti impostati.

    \subsection{Limitazioni delle fonti dati}

    La preoccupazione principale del team si pone sull'assenza di molti valori da un gran numero di cibi confezionati, questo è dovuto in maggior parte alla legge, in quanto le aziende che producono determinati alimenti non sono obbligate a svolgere tutte le analisi (soprattutto dei micronutrienti) necessarie per avere un elenco comprensivo di tutti i valori nutrizionali dell'alimento.

    Questo problema è purtroppo oltre lo scopo del progetto, siccome per avere veramente tutti i dati dei prodotti, bisognerebbe spostare la responsabilità di svolgere le analisi dalle aziende produttrici agli sviluppatori dell'applicazione, per creare un database riservato a quest'ultima, un obiettivo ben oltre lo scopo di NutriApp.

    \section{Informativa sulla privacy e distribuzione del software}

    \subsection{Informativa sulla privacy}

    All'interno del progetto NutriApp è stato predisposto un documento dedicato al trattamento dei dati personali, redatto in conformità al Regolamento (UE) 2016/679 (GDPR).

    Il documento include informazioni sul trattamento dei dati e come richiedere la loro eliminazione, può essere trovato in forma completa al seguente collegamento:

    \href{https://gitlab.com/ismailbarakat68/nutriapp/-/raw/develop/Documentazione/informativa%20sulla%20privacy.pdf?ref_type=heads&inline=true}{Apri il documento di Informativa sulla Privacy}

    NutriApp raccoglie esclusivamente i dati necessari al funzionamento del servizio. Il trattamento avviene nel rispetto delle normative vigenti in materia di protezione dei dati personali.

    \subsection{Distribuzione del software su piattaforme ufficiali}

    L'applicazione \textbf{Nutriapp} è distribuita ufficialmente tramite le principali piattaforme digitali per dispositivi mobili. La distribuzione attraverso piattaforme ufficiali garantisce maggiore sicurezza, affidabilità e disponibilità degli aggiornamenti.

    \subsubsection*{Piattaforme supportate}

    L'applicazione è disponibile per il download attraverso:

    \begin{itemize}

        \item \textbf{Google Play Store}
        per dispositivi basati su sistema operativo Android.

        \item \textbf{Apple App Store}
        per dispositivi basati su sistema operativo iOS
        (iPhone e iPad).

    \end{itemize}

% ------------------------------------------------------------

    \subsection{Distribuzione del software da terze parti}

    Oltre alla distribuzione tramite piattaforme ufficiali,l'applicazione può essere resa disponibile anche attraverso repository di sviluppo,principalmente per attività di test,debug o installazione manuale.

    \subsubsection*{Repository ufficiale}

    Il pacchetto di installazione dell'applicazione è disponibile nella repository GitLab ufficiale del progetto.

    Il file distribuito è nel formato APK

    \subsubsection*{Avvertenze di sicurezza}

    L'installazione manuale deve essere effettuata esclusivamente utilizzando file provenienti dalla repository ufficiale del progetto. L'utilizzo di file provenienti da fonti non verificate può comportare rischi per la sicurezza del dispositivo e per l'integrità dei dati dell'utente.

    Il team Nutriapp non è da ritenersi responsabile per eventuali danni arrecati agli utenti che hanno installato l'applicazione da una fonte non verificata.

% ------------------------------------------------------------

    \subsection{Recapiti e contatti del team Nutriapp}

    Il team di sviluppo mette a disposizione diversi canali di comunicazione per fornire assistenza tecnica, raccogliere segnalazioni di errore e gestire richieste informative.

    \subsubsection*{Canali di supporto}

    Gli utenti possono contattare il team tramite:

    \begin{itemize}

        \item \textbf{Email di supporto tecnico}
        support@nutriapp.com

        \item \textbf{Repository GitLab ufficiale}
        utilizzata per la segnalazione di bug,
        richieste tecniche e contributi allo sviluppo.

        \item \textbf{Sistema di assistenza integrato}
        disponibile direttamente all'interno dell'applicazione.

    \end{itemize}

    \subsubsection*{Informazioni da fornire in caso di errore}

    Per facilitare l'analisi di eventuali problemi, si consiglia di includere:

    \begin{itemize}

        \item modello del dispositivo utilizzato

        \item versione del sistema operativo

        \item versione dell'applicazione installata

        \item descrizione dettagliata del problema riscontrato

    \end{itemize}

    \subsubsection*{Gestione delle richieste}

    Le richieste inviate tramite i canali ufficiali vengono analizzate dal team tecnico in base alla priorità e alla complessità del problema. Eventuali aggiornamenti critici verranno distribuiti attraverso le piattaforme ufficiali di distribuzione.
\end{document}
